\documentclass[11pt,letterpaper]{article}
\usepackage[margin=1.06in]{geometry}
\usepackage[T1]{fontenc}
\usepackage{lmodern,microtype,amsmath,amssymb,amsthm,mathtools}
\usepackage[strings]{underscore}
\usepackage[hidelinks]{hyperref}
\hypersetup{pdftitle={Chiral transfer and pole forms on smooth homotopy four-spheres},pdfauthor={Edward B. Baker III}}
\setlength{\emergencystretch}{2em}
\numberwithin{equation}{section}
\newtheorem{proposition}{Proposition}[section]
\newtheorem{corollary}[proposition]{Corollary}
\title{Chiral transfer and pole forms\\on smooth homotopy four-spheres}
\author{Edward B. Baker III}
\date{27 November 2026}
\begin{document}
\maketitle

\begin{abstract}
Let \(M\) be a smooth homotopy four-sphere, \(c\) a smooth Riemannian
conformal class and \(p\in M\). We construct all closed self-dual
two-forms on \(M\setminus\{p\}\) with growth at most that of a quartic pole.
Their leading coefficients at \(p\) form a real three-dimensional space,
and each coefficient determines a unique global solution. We identify
this solution operator with the off-diagonal kernel of the chiral transfer
\(d^-A\mapsto d^+A\), including its normalization. A single nowhere-zero
solution gives a symplectic structure on the punctured manifold. Its
asymptotics allow us to prescribe the chart at infinity in a relative
symplectic recognition argument and to extend the resulting diffeomorphism
across the original smooth puncture. Thus smooth Poincar\'e is equivalent
to the existence, on each smooth type, of one class, one pole and one
coefficient direction that works everywhere. We express the remaining
condition on a finite cover and examine what can happen when a direction
first acquires a zero.
\end{abstract}

\section{From isothermal coordinates to a recognition question}

On the Riemann sphere a stereographic coordinate \(z\) has a simple pole.
Its differential is closed, has a double pole, and is nowhere zero on the
complement of that pole. It records a global fact about the coordinate in
a form that can be checked locally. The present approach grew out of
work on generalized isothermal coordinates \cite{BakerCurved2025,BakerIsothermal2026}
and a proposal to use pole forms for smooth recognition
\cite{Baker2026PoleFormsProposal}. There is a second feature of the
two-dimensional construction worth retaining. For a map of the plane
normalized as \(f=z+\text{correction}\), its two complex derivatives
are related by \(f_z=1+Bf_{\bar z}\), with \(B\) the Beurling transform.
The Beltrami equation \(f_{\bar z}=\mu f_z\) then becomes a linear
equation for \(f_{\bar z}\). The transfer between two halves of a
differential is the part of the analogy we use in four dimensions.

Fix a closed, connected, oriented smooth four-manifold \(M\) homotopy
equivalent to \(S^4\). Freedman's theorem identifies its underlying
topological manifold \cite[Theorem~1.6]{Freedman1982}. Smooth
recognition asks for a diffeomorphism with the given smooth structure.
Our object is a real two-form on \(X=M\setminus\{p\}\). It has two
properties with different consequences. Closedness and self-duality
give a linear elliptic problem. Nonvanishing turns a solution into a
symplectic form, since
\[
                    \omega\wedge\omega=|\omega|_g^2\,dV_g .
\]
It is this change from linear analysis to symplectic geometry that
makes a single solution sufficient for recognition.

The first main result, Proposition~\ref{prop:pole}, solves the linear
problem for every smooth conformal class and every marked point.
The second, Proposition~\ref{prop:recognition-pole}, recognizes the
original compact manifold whenever one of those solutions is nowhere
zero. The quantifiers matter. With \(u\) a direction in the coefficient
space, the outstanding assertion is
\begin{equation}\label{eq:openquant}
   \forall M\quad \exists\,c,p,[u]\quad
   \forall q\in M\setminus\{p\},\qquad \omega_u(q)\ne0 .
\end{equation}
The coefficient line may depend on \(M,c,p\). It is fixed as \(q\)
varies. Section~\ref{sec:directions} gives an exact criterion for
\eqref{eq:openquant} on compact pieces of a handle atlas. For a second
connection with generalized isothermal coordinates we examine the
resulting map on twistor directions and identify when it gives a
whole fibre chart.

\paragraph{Evidence and conventions.}
An \emph{import} is used with its hypotheses and source at the point
of application. The mathematical results derived here are project
theorems. A1, A2, the normalization and the marked-end results have
independent written proof replays. The finite-cover criterion is an
elementary consequence derived below. A \emph{certificate} verifies
the specified finite inequalities in the supporting example.
The geometric results await formal verification in Lean.
All conformal classes and forms in the main argument are smooth and
real, and all metrics positive definite.

\section{The chiral transfer on a closed manifold}
\label{sec:transfer}

Normalize the two-form norm by \(|e^i\wedge e^j|=1\) for \(i\ne j\)
in an orthonormal coframe. Under \(\widehat g=e^{2f}g\) the inner
product on \(k\)-forms scales by \(e^{-2kf}\) and volume by \(e^{4f}\).
The defining wedge identity gives
\(
*_{\widehat g}|_{\Lambda^k}=e^{(4-2k)f}*_g.
\)
In degree two both the star and the \(L^2\) inner product are
conformal invariants. The star squares to the identity, splitting
two-forms into two oriented rank-three bundles \(\Lambda^\pm\).
Write \(P^\pm=(1\pm *)/2\) and \(d^\pm=P^\pm d\).

For a smooth one-form \(A\), Stokes's theorem gives
\[
   0=\int_M d(A\wedge dA)=
   \int_M\bigl(|d^+A|^2-|d^-A|^2\bigr)\,dV_g .
\]
Density extends the identity to \(W^{1,2}\). It shows that the exact
two-form \(dA\) is determined by either chiral part. To see that every
\(L^2\) input occurs, use \(D_g=(d_g^*,\sqrt2\,d^-)\) on one-forms.
Integration by parts gives \(D_g^*D_g=\Delta_1\) and
\(D_gD_g^*=\Delta_0\oplus\Delta_2^-\). The first kernel vanishes since
\(H^1(M;\mathbb R)=0\). The second consists of constant functions in
the scalar slot since \(H^2(M;\mathbb R)=0\).
Compact Hodge theory therefore solves
\(D_g A=(0,\sqrt2\,\alpha)\) uniquely for every
\(\alpha\in L^2(\Lambda^-)\), with \(A\in W^{1,2}\). Reversal of
orientation gives the opposite construction. Thus
\begin{equation}\label{eq:transfer}
 T_c:L^2(\Lambda^-)\longrightarrow L^2(\Lambda^+),\qquad
 T_c(d^-A)=d^+A
\end{equation}
is a surjective isometry. The potential in Coulomb gauge uses a
representative of \(c\); the graph of exact forms and hence \(T_c\)
depend only on \(c\). For \(F=dA\) in real abelian curvature
normalization, \(F^+=T_c F^-\). Donaldson and Sullivan give the
inverse arrow \cite[\S2(ii), (2.9)--(2.13)]{DonaldsonSullivan1989}.
For the Euclidean operator on differential forms see Iwaniec and
Martin \cite[\S8]{IwaniecMartin1993}.

At \(p\), a concentrated anti-self-dual source should therefore generate
a self-dual field elsewhere. We first classify fields by growth, without
choosing a normalization for a point source. Set
\begin{equation}\label{eq:pole-space}
 \mathcal H_p(c)=\bigl\{\omega\in\Omega^2(X):
   d\omega=0,\quad *_c\omega=\omega,\quad
   |\omega|_g=O(d_g(p,\cdot)^{-4})\bigr\}.
\end{equation}
Smooth conformal representatives have locally comparable norms and
distances, so the growth condition is intrinsic. On the round sphere
every constant self-dual form in oriented stereographic coordinates
belongs to this space. Inversion of the coordinates accounts for the
power four: each covector contributes a factor of inverse distance
squared. The analytic question is whether every such constant at the
marked end extends, uniquely, on an arbitrary \(X\).

\section{Existence and classification of pole forms}\label{sec:pole-existence}

Fix \(g\in c\) and oriented orthonormal \(g\)-normal coordinates \(y\) at
\(p\). Put \(\rho=|y|\), \(x=I(y)=y/\rho^2\), and \(r=|x|\).
For a coefficient array in these Cartesian charts,
\(O_\infty(r^\gamma)\) means that each derivative of order \(j\) is
\(O(r^{\gamma-j})\). We use the same convention with \(\rho\) near
the puncture. Let \(R_y=\mathrm{Id}-2\hat y\otimes\hat y\), with
\(\hat y=y/|y|\).

\begin{proposition}[A1, replayed project theorem]\label{prop:pole}
For every smooth conformal class \(c\) and every \(p\in M\), the space
\(\mathcal H_p(c)\) of \eqref{eq:pole-space} is three-dimensional.
Each \(\eta\in\Lambda_p^-\) determines a unique member \(\omega_\eta\)
with expansion
\[
 \omega_\eta=\rho^{-4}R_y^*\eta+O_\infty(\rho^{-2}),
\]
where the coefficients of \(\eta\) are held constant in the normal
chart. The assignment \(\eta\mapsto\omega_\eta\) is a linear
isomorphism. On the inverted end there is a one-form
\(a=O_\infty(r^{-1})\) such that \(\omega_\eta=\omega^0+da\), where
\(\omega^0\) is the constant two-form with the coefficients of \(\eta\)
in the \(x\)-chart.
\end{proposition}

\begin{proof}
Choose a smooth positive conformal multiple \(h\) of \(g\) on \(X\)
equal to \(\rho^{-4}g\) near \(p\). From \(D_xy=r^{-2}R_x\) we obtain
\[
 h(x)=R_x^{\mathsf T}g(x/r^2)R_x=\mathrm{Id}+O_\infty(r^{-2}).
\]
In normal coordinates \(g-\mathrm{Id}\) vanishes to second order
in \(y\). Since \(y=x/r^2\) is homogeneous of degree \(-1\) and \(R_x\)
is homogeneous of degree zero, each Cartesian derivative in \(x\)
improves the decay by one order. The puncture supplies the only end of
\(X\), and uniform comparability with the Euclidean metric there gives
completeness of \(h\). The inversion reverses orientation, so
Euclidean self-duality for the orientation of \(X\) is anti-self-duality
for the usual coordinate orientation of \(x\). The stars of \(h\) and
\(g\) agree on two-forms, and near \(p\) their two-form norms satisfy
\(|\omega|_h=\rho^4|\omega|_g\).

\paragraph{The weighted operator.}
Let \(P_h^-=(1-*_h)/2\) on two-forms, let \(d_h^*\) be the formal
adjoint of \(d\), and write \(\Delta_h=dd_h^*+d_h^*d\). Set
\[
 D=(d_h^*,\sqrt{2}\,P_h^-d),\qquad
 D^*(f,\beta)=df+\sqrt{2}\,d_h^*\beta
 \quad(\beta\in\Lambda_h^-).
\]
The identities \(D^*D=\Delta_h\) on one-forms and
\(DD^*=\Delta_h\) on functions and anti-self-dual two-forms give
square ellipticity. The first uses \(d_h^**_hd=0\). For the second,
\(d_h^*\beta=*_h d\beta\) implies
\(d_h^*d\beta=-*_h dd_h^*\beta\), so
\(2P_h^-dd_h^*\beta=\Delta_h\beta\). The off-diagonal terms vanish
by \(d^2=(d_h^*)^2=0\).

Use \(W^{k,2}_\delta\) with the usual Sobolev norm on the compact part
and squared norm on the end equivalent to
\[
 \sum_{j=0}^k\int r^{2(j-\delta)}|\partial^j u|^2\,r^{-4}\,d^4x,
 \qquad L^2_\delta=W^{0,2}_\delta.
\]
For large \(r\), the restriction of \(P_h^-\) identifies the flat
anti-self-dual bundle with \(\Lambda_h^-\), using the same orientation
for both. In these end frames the flat operator is \(D_0\).
The first- and zeroth-order coefficients of \(D-D_0\) have orders
\(O_\infty(r^{-2})\) and \(O_\infty(r^{-3})\), respectively.
The same bounds hold for \(D^*-D_0^*\) and for
\(Q-Q_0\), where \(Q=d+d_h^*\) and \(Q_0=d+d_0^*\) act on all forms.

To apply the cylindrical Fredholm theorem, set \(z=\log(r/R_0)\) on
the end and \(A=rD\), extending \(r\) smoothly and positively on the
compact part. In cylindrical coordinates with Cartesian component
frames, \(rD_0\) is translation invariant, and the coefficients of
\(rD-rD_0\) and all their cylindrical derivatives are \(O(r^{-2})\).
A global comparison operator is obtained by interpolating between
\(A\) on the compact part and \(rD_0\) on the end. On a sufficiently
distant interpolation annulus the principal symbols are uniformly close
on unit cylindrical covectors, so the comparison operator remains
elliptic. The end measure is \(r^{-4}d^4x=dz\,dS\) and the exponential
factor in the cylindrical norm is \(e^{-\delta z}\), up to a constant.
Multiplying the target by \(r\) shifts its weight from \(\delta-1\)
to \(\delta\). Thus the weighted Fredholm theorem for a smooth square
elliptic system with coefficients converging to a translation-invariant
elliptic system on a cylindrical end applies to
\[
 D\colon W^{1,2}_\delta\longrightarrow L^2_{\delta-1}.
\]
In the notation of the theorem the exceptional cylindrical weights are
\(\operatorname{Im}\lambda\) for nonzero pencil modes
\(e^{i\lambda z}v(\theta)\) of the limit operator. With \(q=i\lambda\),
these are the Cartesian weights \(\delta=\operatorname{Re}q\) for
which \(D_0(r^qv)=0\) with \(v\) smooth on \(S^3\).
Generalized modes with polynomial factors in \(z\) introduce no new
weights, since their highest coefficient solves the ordinary pencil
equation. The index and nullspaces are constant between consecutive
exceptional weights
\cite[Theorem~6.2 and Lemma~7.1]{LockhartMcOwen1985Elliptic}.

Since \(D_0^*D_0=\Delta_0\), every Cartesian coefficient of
\(r^qv\) is harmonic. The eigenvalues \(\ell(\ell+2)\) on the unit
\(S^3\) give \(q(q+2)=\ell(\ell+2)\) for \(\ell=0,1,\ldots\).
They follow from restricting homogeneous harmonic polynomials to the
sphere, whose restrictions span by the harmonic decomposition of
polynomials. In \([-3,0]\) the candidates are \(-3,-2,0\).
At \(-2\), the coefficient vector has the form \(r^{-2}v\) with
\(v\) constant, so its codifferential is \(2r^{-4}x\cdot v\).
This vanishes identically only for \(v=0\).
The one-form \(d(r^{-2})\) realizes \(-3\), and constant one-forms
realize \(0\). These are the only exceptional weights in the interval.

Local elliptic regularity gives smoothness for distributional solutions
with smooth data. To track decay, rescale \(x=Rs\) on fixed nested
annuli \(B\Subset B'\). For \(P=D,D^*,Q\) or their flat versions, the
operator \(P_R\) induced by \(RP\) has coefficients converging in every
\(C^m\) norm on \(B'\) to those of the corresponding flat operator in
\(s\). The interior parametrix estimates
\cite[estimate~(2.4)]{LockhartMcOwen1985Elliptic} take the uniform form
\[
 \|u_R\|_{H^{m+1}(B)}
 \le C_m\bigl(\|R(Pu)_R\|_{H^m(B')}+\|u_R\|_{L^2(B')}\bigr),
 \qquad u_R(s)=u(Rs).
\]
Uniformity on the end also follows by absorbing the small coefficient
error in the flat estimates. The rescaled \(L^2\) term for \(u\in
L^2_\delta\) is bounded by \(C R^\delta\|u\|_{L^2_\delta}\), and a
pointwise \(O(r^\delta)\) bound gives the same power of \(R\).
If \(Pu=O_\infty(r^\gamma)\), the rescaled source term is
\(O(R^{\gamma+1})\). Sobolev embedding gives derivative bounds
\(R^j|\partial^j u|\le C_j(R^\delta+R^{\gamma+1})\) on the smaller
annulus, with only the first term for a homogeneous solution.

\paragraph{The kernel and cokernel.}
Work first at \(\delta=-3/2\). If \(Da=0\), regularity gives
\(a=O_\infty(r^{-3/2})\), and \(da\) is self-dual. For the
truncation \(X_R\) at \(r=R\), Stokes gives
\[
 \int_{X_R}|da|_h^2\,dV_h
 =\int_{\partial X_R}a\wedge da=O(R^{-1}),
\]
using boundary area \(O(R^3)\) and coefficient orders \(-3/2\)
and \(-5/2\). Hence \(da=0\).
By excision \(H^2(M,X;\mathbb R)\) is the local cohomology group
\(H^2(B^4,B^4\setminus\{0\};\mathbb R)=0\). As \(H^1(M;\mathbb R)=0\),
the relative sequence gives \(H^1(X;\mathbb R)=0\). Thus \(a=dF\)
with \(\Delta_hF=0\).
Integrating \(dF\) along rays shows that \(F\) converges uniformly
along them with tail \(O(r^{-1/2})\). The oscillation on each sphere is
also \(O(r^{-1/2})\), since the sphere diameters grow linearly. Thus
\(F\) has a common limit at infinity. The maximum principle on \(X_R\)
gives \(F\) constant.

The \(h\)-volume pairing identifies the dual of \(L^2_{\delta-1}\)
with \(L^2_{-3-\delta}\). Indeed the two weight exponents sum to
\(-4\), so the product of the norm factors \(r^{-\delta+1}\) and
\(r^{3+\delta}\) converts \(r^{-4}d^4x\) to \(d^4x\). The density
and pointwise inner product for \(h\) are uniformly equivalent to
their Euclidean versions. A cokernel representative at \(\delta=-3/2\)
therefore solves
\(D^*(f,\beta)=0\) distributionally in \(L^2_{-3/2}\).
It is smooth and \(O_\infty(r^{-3/2})\).
Applying \(d_h^*\) gives \(\Delta_h f=0\), and the maximum principle
with \(f\to0\) gives \(f=0\). Then \(d_h^*\beta=0\) and chirality
implies \(d\beta=0\). In the compact metric,
\(|\beta|_g=O(\rho^{-5/2})\).

For either chirality, a closed two-form with compact-metric norm
\(O(\rho^{-b})\) for \(b<3\) has a removable singularity. The form
is locally integrable, and testing its exterior derivative with a cutoff
that changes from zero to one on \(\epsilon<\rho<2\epsilon\) costs
\(O(\epsilon^{3-b})\). The factor \(\epsilon^{-1}\) from the cutoff
derivative is multiplied by \(O(\epsilon^{-b})\) from the form and
\(O(\epsilon^4)\) from the annular volume. The extension is therefore
closed distributionally across \(p\). The chiral identity gives distributional coclosedness, and
elliptic regularity for \(d+d_g^*\) gives a smooth extension. Since
\(H^2(M;\mathbb R)=0\), the extension is exact, and Stokes and
\(\beta\wedge\beta=\pm|\beta|_g^2dV_g\) make it zero. This applies to
the cokernel form with \(b=5/2\).

Fredholmness now gives the inverse at \(-3/2\). Nullspace and index
constancy give an inverse for every \(-3<\delta<0\). The same inverse
maps \(W^{k,2}_{\delta-1}\) boundedly to \(W^{k+1,2}_\delta\) for
each integer \(k\ge0\). Indeed the local rescaled bounds, summed in
weighted squares over dyadic annuli of finite overlap and combined
with compact interior estimates, bound the higher domain norm by the
higher data norm and the \(L^2_\delta\) norm of the solution. The
order-one inverse controls this last term.

\paragraph{The endpoint estimates.}
For \(P_0=D_0\) or \(Q_0\), let
\(E=P_0^*\Gamma\), where \(\Gamma(x)=(4\pi^2|x|^2)^{-1}\)
is the four-dimensional fundamental solution for the positive
Laplacian. It has this normalization by flux through a sphere.
Since \(P_0P_0^*=\Delta_0\), \(P_0E=\delta_0\mathrm{Id}\)
and \(|E(x)|\le C|x|^{-3}\).
If \(G\) is smooth on \(\mathbb R^4\) and \(G=O(r^{-k})\) for
\(k=2\) or \(3\), then \(E*G\) converges absolutely and has
order \(O(r^{1-k})\). Split the convolution integral at
\(|w|=r/2\) and \(2r\). The respective bounds are
\[
 Cr^{-3}\int_{|w|\le r/2}(1+|w|)^{-k}d^4w,qquad
 Cr^{-k}\int_{|x-w|\le3r}|x-w|^{-3}d^4w,qquad
 C\int_{2r}^\infty t^{-k}dt.
\]
Each is \(O(r^{1-k})\). Local integrability and the same tail bound
justify testing the convolution against compactly supported functions,
so \(P_0(E*G)=G\) in distributions. If a bounded section \(u\) solves
\(P_0u=G\) globally, its difference from \(E*G\) is entire
coefficientwise harmonic and hence constant. This constant is zero
when \(u\to0\).

\paragraph{Existence and exhaustion.}
For the prescribed \(\eta\), let \(\omega^0\) be its constant form
on the end, Euclidean self-dual for the orientation of \(X\).
Use the radial vector field \(\mathcal R=\sum_j x^j\partial_{x^j}\).
Cut off \(A^0=\tfrac12\iota_{\mathcal R}\omega^0\) so that it extends
smoothly by zero on the compact part and \(dA^0=\omega^0\) far out.
Then \(P_h^-dA^0=O_\infty(r^{-2})\). Solve
\[
 Da=(0,-\sqrt{2}\,P_h^-dA^0)
\]
at any \(-1<\delta<0\), since the source lies in
\(L^2_{\delta-1}\). The rescaled estimates give
\(a=O_\infty(r^\delta)\). Thus \(D_0a=O_\infty(r^{-2})\), since
the additional error \((D_0-D)a\) is \(O_\infty(r^{\delta-3})\).
Cut off \(a\) in the flat end frames and extend by zero to
\(\mathbb R^4\). The preceding convolution estimate with \(k=2\)
applies, and the entire harmonic residual vanishes because the cutoff
section tends to zero. This gives \(a=O(r^{-1})\), which the flat annular estimates
upgrade to \(O_\infty(r^{-1})\). The two-form \(d(A^0+a)\) is closed by
construction and self-dual by the second component of the equation
for \(a\). It equals \(\omega^0+da\) far out with
\(da=O_\infty(r^{-2})\). The convolution is what gives the stated
endpoint, since the weighted step gives only
\(da=O_\infty(r^{\delta-1})\).

Conversely, every \(\omega\in\mathcal H_p(c)\) is bounded in the
Cartesian end frames by the compact-to-end norm relation. Closedness
and self-duality imply \(Q\omega=0\). The rescaled estimates give
\(\omega=O_\infty(1)\), so \(Q_0\omega=O_\infty(r^{-3})\).
Cut off \(\omega\) on the end and apply the convolution estimate
with \(k=3\). The constant harmonic residual is a two-form because
the cutoff section is a two-form and the convolution tends to zero.
Hence \(\omega=\omega^0+O(r^{-2})\) for a constant two-form.
Passing to the limit in \(P_h^-\omega=0\) gives
\(P_0^-\omega^0=0\), and flat interior estimates give the differentiated
remainder. Pullback under \(I\), with \(D_yI=\rho^{-2}R_y\),
takes \(\omega^0\) to \(\rho^{-4}R_y^*\eta\) and takes
\(O_\infty(r^{-2})\) to \(O_\infty(\rho^{-2})\).
Subtracting the constructed solution for this \(\eta\) leaves a
removable closed self-dual form with \(b=2\), which vanishes by the
removability argument above. The construction and uniqueness are linear,
and the constant chiral space has dimension three.
\end{proof}

The order of the remainder can be attained in normal coordinates.
On \(S^4\), take \(g\) equal near \(p\) to the product of a spherical
surface of curvature one and a flat plane, and interpolate with a
smooth metric away from \(p\). On the ray \(y=(\rho,0,0,0)\) in product
normal coordinates, the spherical polar metric gives
\(g=\mathrm{diag}(1,b^2,1,1)\), where \(b=\sin\rho/\rho\).
Reflection in the inversion formula preserves this diagonal matrix,
so \(h\) has the same entries on the corresponding ray. If
\(dx^{ij}=dx^i\wedge dx^j\) and \(\omega^0=dx^{12}-dx^{34}\), the star
of the reversed orientation gives
\[
 (*_h-1)\omega^0
 =(b-1)dx^{12}+(1-b^{-1})dx^{34}
 =-\frac{r^{-2}}6(dx^{12}+dx^{34})+O(r^{-4}).
\]
If the pole form with this coefficient had error \(o(r^{-2})\) on the
ray, its self-duality and boundedness of \(*_h\) would force
\((*_h-1)\omega^0=o(r^{-2})\), contrary to the displayed coefficient.

This completes the growth classification. To identify these solutions as the
columns of \eqref{eq:transfer}, it remains to compute the normalization
of a point source.

\section{The transfer kernel and its pole coefficient}
\label{sec:transfer-kernel}

Fix a representative \(g\in c\) and let \(dV_g\) denote its positive
volume density. We write \(K_g(q,p)\) for the kernel of \eqref{eq:transfer}
relative to this density at the source \(p\). Thus for a smooth input
\(\alpha\) whose support omits \(q\),
\[
 (T_c\alpha)(q)=\int_M K_g(q,p)\alpha(p)\,dV_g(p).
\]
For distinct \(p,q\), the map \(K_g(q,p)\) takes \(\Lambda^-_p\) to
\(\Lambda^+_q\). Write \(P_g(q,p)\eta=\omega_\eta(q)\) for the solution
of Proposition~\ref{prop:pole} with leading coefficient
\(\eta\in\Lambda^-_p\). The relation is
\begin{equation}\label{eq:kernel-pole-normalization}
 K_g(q,p)=\frac{2}{\pi^2}P_g(q,p),\qquad q\ne p.
\end{equation}

To see the smoothness needed for the pointwise kernel, let
\(\Delta_g=dd_g^*+d_g^*d\) be the Hodge Laplacian on two-forms and
\(\pi_g^+\) the self-dual projection. As \(H^2(M;\mathbb R)=0\),
\(dd_g^*\Delta_g^{-1}\) is the orthogonal projection onto exact
two-forms. The exact graph gives
\(T_c=2\pi_g^+dd_g^*\Delta_g^{-1}|_{\Lambda^-}\).
Indeed, the projection of \(\alpha\in\Lambda^-\) onto the graph is the
exact form \((\alpha+T_c\alpha)/2\), since the difference is
orthogonal to the graph by the isometry of \(T_c\). It follows that
\(\Delta_g T_c\) and \(T_c\Delta_g\) both equal the differential operator
\(2\pi_g^+dd_g^*|_{\Lambda^-}\). Off the diagonal the distributional
kernel is therefore annihilated by the Hodge Laplacians acting at the
target and, by formal adjunction, at the source. Their sum is elliptic on
the product, which gives joint smoothness there.

For the normalization, use the oriented orthonormal normal coordinates
\(y\) at \(p\) of Proposition~\ref{prop:pole}. Set \(\rho=|y|\),
\(n=y/\rho\), and \(I(y)=y/\rho^2\), and identify \(\eta\) with its
constant-coefficient extension in the chart. The locally integrable
one-form \(b=\tfrac12\rho^{-4}\iota_y\eta\) satisfies
\(db=I^*\eta\) for \(y\ne0\), since it is the pullback under \(I\) of the
linear primitive of the constant form. The angular mean of the leading term
\(I^*\eta=\rho^{-4}(\eta-2n^\flat\wedge\iota_n\eta)\) vanishes.
Indeed \(\int_{S^3}n_i n_j\,dS=(2\pi^2/4)\delta_{ij}\), and
\(\sum_i e^i\wedge\iota_{e_i}\eta=2\eta\) for a Euclidean orthonormal
frame. The spherical principal value thus exists. Differentiating
\(\mathbf{1}_{\{\rho>\varepsilon\}}b\) gives a boundary term with
positive radial conormal \(d\rho\). Its limit is a point mass with
coefficient
\[
 \frac12\int_{S^3}n^\flat\wedge\iota_n\eta\,dS
 =\frac{\pi^2}{2}\eta .
\]
In coefficient distributions in the normal chart this reads
\[
 db=\operatorname{pv}(I^*\eta)+\frac{\pi^2}{2}\eta\,\delta_0 .
\]

Let \(\delta_p^g\) be the scalar point mass relative to \(dV_g\), so
that pairing \(\eta\delta_p^g\) with a smooth two-form \(\gamma\)
gives \(\langle\eta,\gamma(p)\rangle_g\). The normal-coordinate volume
density at \(p\) is one. The remainder \(\omega_\eta-I^*\eta\) is closed
on the deleted ball and is \(O(\rho^{-2})\). Its locally integrable
extension is closed across \(p\), since the boundary integral against a
smooth test one-form is \(O(\rho^3\rho^{-2})\). Consequently the
two-form distribution
\[
 F=\operatorname{pv}(\omega_\eta)+\frac{\pi^2}{2}\eta\delta_p^g
\]
is closed on \(M\). Its principal-value term is self-dual, as follows by
testing self-duality on every deleted ball and passing to the limit.

Reverse the orientation \(o\) to exchange the two chiral projections.
The graph in \eqref{eq:transfer} then gives
\(T_c^{-o}=(T_c^o)^{-1}=(T_c^o)^*\).
For a smooth self-dual test form \(\beta\), take a smooth potential \(B\)
with \(d^+B=\beta\), using the compact inverse on a homology
four-sphere \cite[\S2(ii), (2.9)--(2.13)]{DonaldsonSullivan1989}.
Then \(d^-B=(T_c^o)^*\beta\). Distributional Stokes gives
\[
 0=\int_M F\wedge dB
  =\big\langle\operatorname{pv}(\omega_\eta),\beta\big\rangle_{L^2(g)}
   -\frac{\pi^2}{2}\big\langle\eta,[(T_c^o)^*\beta](p)\big\rangle_g .
\]
The minus sign is the wedge pairing on anti-self-dual forms. Thus
\((2/\pi^2)\operatorname{pv}(\omega_\eta)\) is the distributional output
of \(T_c^o\) on \(\eta\delta_p^g\). For test forms supported away from
\(p\), this output is represented by the smooth off-diagonal kernel
column. This proves \eqref{eq:kernel-pole-normalization}, with the
positive constant for the arrow \(d^-A\mapsto d^+A\).

The conformal class fixes both \(T_c\) and the chiral bundles. Its
density-valued kernel \(K_g(q,p)dV_g(p)\) is therefore independent of
\(g\). In particular, for \(\widehat g=e^{2f}g\),
\[
 K_{\widehat g}(q,p)=e^{-4f(p)}K_g(q,p)
\]
as maps on the same covariant source forms. If \(\eta_g\) and
\(\eta_{\widehat g}\) label a fixed output two-form under
\(P_g\) and \(P_{\widehat g}\), respectively,
\eqref{eq:kernel-pole-normalization} and uniqueness of the coefficient
give \(\eta_{\widehat g}=e^{4f(p)}\eta_g\). Thus the intrinsic label is
\(\eta_g\otimes dV_g(p)^{-1}\). In aligned orthonormal coframes the
components of that fixed output's coefficient acquire a factor
\(e^{2f(p)}\) because each coframe covector scales by \(e^{f(p)}\).

The change of normal coordinates gives the same weight for a
nonconstant conformal factor. Let \(z\) be the new coordinate with
aligned frames, \(a=e^{f(p)}\), and \(v=\nabla_g f(p)\) in the old
frame. At \(p\) the difference of the two connections on \((u,u)\) is
\(2(v\cdot u)u-|u|^2v\). The geodesic equation and inversion of its
Taylor expansion give
\[
 z=a\bigl(y+(v\cdot y)y-\tfrac12|y|^2v+O(|y|^3)\bigr).
\]
The squared radius is
\(a^2|y|^2(1+v\cdot y+O(|y|^2))\). Setting \(x=I(y)\) consequently
gives \(I(z)=a^{-1}(x-v/2+O(|x|^{-1}))\). The differentiated Taylor
remainder gives \(D_xI(z)=a^{-1}(\mathrm{Id}+O(|x|^{-2}))\).
The translation leaves the constant two-form unchanged. The linear
factor contributes twice to its pullback, and the residual term has
the decay of the remainder in the inverted chart.

Finally, taking the \(L^2\) adjoint exchanges source and target in the
kernel. Since \(T_c^{-o}=(T_c^o)^*\), this gives
\begin{equation}\label{eq:kernel-orientation-reciprocity}
 K_g^{-o}(p,q)=K_g^o(q,p)^{*_{g(p),g(q)}},\qquad p\ne q,
\end{equation}
where the adjoint uses the positive fibre inner products at the two
points. In orientation \(-o\), the source at \(q\) is
\(\Lambda^+_{o,q}\) and the target at \(p\) is \(\Lambda^-_{o,p}\).

\section{Recognition from a nowhere-zero pole form}

\begin{proposition}[A2, replayed project theorem]\label{prop:recognition-pole}
Let \(M\) be a closed, connected, oriented smooth homotopy four-sphere
with a smooth Riemannian conformal class \(c\) and a point \(p\).
If some \(\omega\in\mathcal H_p(c)\) is nowhere zero on
\(X=M\setminus\{p\}\), then \(M\) is diffeomorphic to the standard
smooth \(S^4\).
\end{proposition}

\begin{proof}
For \(g\in c\), self-duality gives
\(\omega\wedge\omega=|\omega|_g^2\,dV_g>0\) on \(X\).
Thus the closed form \(\omega\) is symplectic and induces the given
orientation. In the inverted normal chart \(x=y/|y|^2\), write
\(\omega=\beta+da\), with \(|a|=O(r^{-1})\) and \(|da|=O(r^{-2})\)
in the Euclidean norm, where \(r=|x|\). These are the end estimates
from Proposition~\ref{prop:pole} and its proof. The constant
\(\beta\) is nonzero by the leading-coefficient isomorphism.
It is Euclidean self-dual in the inherited orientation, which is
opposite to the usual orientation of the \(x\)-coordinates.

Take a smooth cutoff \(\chi_R\) equal to one for \(r\le R\) and zero
for \(r\ge 2R\), with \(|d\chi_R|\le C/R\). On the end put
\(\Omega_X=\beta+d(\chi_Ra)\) and retain \(\omega\) on the core.
The expressions agree on an open overlap and define a closed smooth
form. On the transition annulus,
\[
 \Omega_X-\beta=\chi_R\,da+d\chi_R\wedge a=O(R^{-2}).
\]
In a Euclidean orthonormal coframe for the inherited orientation,
\(\beta=b(e^1\wedge e^2+e^3\wedge e^4)\), with
\(b=|\beta|/\sqrt2>0\). As a skew map it has least singular value
\(b\). Take \(R\) large enough that the annular error has operator
norm less than \(b\). The segment from \(\beta\) to \(\Omega_X\)
then consists of nondegenerate forms at each annular point and
preserves the orientation. Hence \(\Omega_X\) is symplectic everywhere
and is exactly \(\beta\) for \(r\ge 2R\).

Let \(\omega_{\mathrm{std}}=dz^1\wedge dz^2+dz^3\wedge dz^4\).
The coframe above gives an orthogonal \(Q\) with
\(Q^*\omega_{\mathrm{std}}=\beta/b\), where \(\det Q=-1\) in
the usual \(x\)-orientation. The map \(e\colon x\mapsto z=sQx\),
where \(s=\sqrt b\), is therefore a symplectic end chart for
\(\Omega_X\). As a similarity it sends round exteriors to round
exteriors. Fix a round truncation sphere in this chart and let \(K\)
be its compact core in \(X\). Then \(X\) deformation retracts onto
\(K\) by radial retraction on the end. Excision gives
\(H_i(M,X;\mathbb Z)\cong H_i(B^4,B^4\setminus\{0\};\mathbb Z)\),
concentrated in degree four, so the exact sequence of the pair and
\(H_2(M;\mathbb Z)=0\) imply \(H_2(K;\mathbb Z)=0\).

Choose \(L\) large enough that the truncating ball and an open
annulus around its boundary lie inside the polydisc
\(P_L=D^2(L)\times D^2(L)\). In
\(Y_0=(S^2\times S^2,\sigma\oplus\sigma)\) take the same round
area form \(\sigma\) on each factor, scaled to total area
\(A=\pi L^2\). Let
\(D=S_1\cup S_2\), where \(S_1=S^2\times\{q\}\) and
\(S_2=\{q\}\times S^2\). Area coordinates on the two factors identify
\(Y_0\setminus D\) with \((P_L,\omega_{\mathrm{std}})\). In a
stereographic coordinate \(\zeta\) on \(S^2\setminus\{q\}\),
the disk map \(\zeta\mapsto L\zeta/\sqrt{1+|\zeta|^2}\)
has area Jacobian \(L^2(1+|\zeta|^2)^{-2}\), the density of
\(\sigma\) in that chart.

Replace the truncating ball in this product with \(K\), identifying
collars by \(e\). The collars fill opposite sides of the seam and
the forms on them are both \(\omega_{\mathrm{std}}\) in the same
coordinates. The result is a closed symplectic manifold
\((Y,\Omega_Y)\) with its prescribed split neighborhood of \(D\).
The spheres there have trivial normal bundles, area \(A\), and
one transverse intersection of sign \(+1\) in their symplectic
orientations. Since the seam is a three-sphere, Mayer--Vietoris
in degree two identifies \(H_2(Y;\mathbb Z)\) with
\(H_2(Y_0;\mathbb Z)\), using \(H_2(K;\mathbb Z)=0\) and the
fact that removing a ball from \(Y_0\) leaves its second homology
unchanged. The two divisor classes are thus an integral basis.
A smooth map from a sphere with image class \(m[S_1]+n[S_2]\) has
\(\Omega_Y\)-period \((m+n)A\).

Gromov's criterion and relative extension
\cite[\S2.4, product criterion and following remark, p.~337]{Gromov1985PseudoHolomorphic}
apply to this sphere pair. The symplectomorphism
\(Y\to Y_0\) extends the prescribed product identification on a
neighborhood of \(D\). After removing \(D\), its target is
\(P_L\), and its domain is the portion of \(X\) inside this
polydisc in the capped description. The complement of the
neighborhood of agreement in \(Y_0\) is compact and contained
in \(P_L\). Choose \(\ell<L\) so that \(P_\ell\) contains
that complement as well as the truncating ball and its collar.
On \(P_L\setminus\overline{P_\ell}\) the relative map is
exactly the prescribed chart \(e\). The map and its inverse
therefore paste with \(e\) and \(e^{-1}\) across this open region
to give a global symplectomorphism
\(F\colon(X,\Omega_X)\to(\mathbb R^4,\omega_{\mathrm{std}})\).
Since \(P_L\) is bounded, \(F=sQx\) on a whole Euclidean tail.

Use \(w=z/|z|^2\) as the smooth target coordinate at infinity
of the standard \(S^4\). In the original smooth chart \(y\) near
\(p\), the map is exactly
\[
 w(F(y))=\frac{sQ(y/|y|^2)}{|sQ(y/|y|^2)|^2}
       =s^{-1}Qy .
\]
It extends across \(p\) with an invertible derivative and a
smooth linear inverse in these charts. This extends \(F\)
to the desired diffeomorphism.
\end{proof}

\begin{corollary}[Reformulation]\label{cor:equivalence}
Every smooth homotopy four-sphere is diffeomorphic to \(S^4\) if
and only if each such \(M\) admits a smooth Riemannian conformal
class \(c\) and a point \(p\) for which one element of
\(\mathcal H_p(c)\) is nowhere zero on \(M\setminus\{p\}\).
\end{corollary}

\begin{proof}
Proposition~\ref{prop:recognition-pole} gives one implication.
For the other, take an orientation-preserving diffeomorphism to
the standard sphere, composing with a round reflection if necessary.
The round class has conformally Euclidean stereographic coordinates
away from the pole. A nonzero constant self-dual two-form in coordinates
oriented as the punctured sphere is closed and nowhere zero. Its
coefficients in the inverted smooth chart are \(O(|w|^{-4})\),
as required for membership in \(\mathcal H_p(c)\). Pull back the
round class, the pole and the form by the diffeomorphism.
\end{proof}

In the existence statement the class and the pole can depend on
the smooth sphere. For any chosen \((c,p)\), A1 identifies a single
pole form with a fixed leading coefficient. Rescaling a nonzero
coefficient to unit norm gives the equivalent requirement of one
unit direction whose form is nonvanishing throughout \(X\).

\section{One coefficient direction on a finite cover}
\label{sec:directions}

Fix \((c,p)\) and a representative \(g\). Put \(V=\Lambda^-_p\),
with its \(g(p)\) inner product. The three-dimensional space furnished
by A1 gives an evaluation map at every point,
\[
      L_q:V\longrightarrow\Lambda^+_q,\qquad L_qu=\omega_u(q).
\]
Changing \(u\) by a nonzero scalar changes no zero set. Thus the
natural space of choices for recognition is \(\mathbb P(V)\simeq
\mathbb{RP}^2\). At a point with \(\operatorname{rank} L_q=2\), one
projective direction vanishes; at rank one, a whole projective line
vanishes. A rank-zero point excludes every direction for this
particular class and pole.

\begin{proposition}[A finite-cover criterion]\label{prop:cover}
Choose the asymptotically Euclidean representative \(h\) of A1.
There is a compact core \(C\subset X\) outside which \(L_q\) is
invertible. Cover \(C\) by finitely many compact pieces \(C_j\) in
smooth coordinate charts. For a unit \(u\in V\), define
\[
 m_j(u)=\min_{q\in C_j}|L_qu|_h,\qquad
 m(u)=\min_j m_j(u).
\]
There is a zero-free member of \(\mathcal H_p(c)\) if and only if
\(\max_{|u|=1}m(u)>0\). Equivalently the compact set
\[
  \mathcal B(c,p)=
  \bigcup_{q\in C}\mathbb P(\ker L_q)
  \ \subseteq\ \mathbb P(V)
\]
is a proper subset, with \(\mathbb P(0)=\varnothing\).
\end{proposition}

\begin{proof}
Apply the \(O(r^{-2})\) end estimate to an orthonormal basis of \(V\).
In the Cartesian end frame the limit of \(u\mapsto L_qu\) is the
inclusion of \(V\) as constant two-forms, an isometry for the flat end
norm. The errors on a basis bound the operator norm of the error for
every unit \(u\) at once. As \(h\) tends to the Euclidean metric, we
may take \(C\) so that \(|L_qu|_h\ge 1/2\) outside \(C\) for every
unit \(u\). Each \(m_j\) is continuous: its change between \(u\) and
\(v\) is bounded by \(\sup_{C_j}\|L_q\|\,|u-v|\).
The maximum over the coefficient sphere is therefore attained. A
positive maximum supplies one direction with a positive minimum on
all the pieces, and the exterior estimate completes the argument. In
the opposite direction a nonvanishing form has a positive minimum on
\(C\) by compactness. Finally the incidence relation \(L_qu=0\) is
closed on \(C\times\mathbb P(V)\), so its projection \(\mathcal B(c,p)\)
is compact. Its complement consists exactly of the desired lines.
\end{proof}

The order of the extrema is important. The quantity
\(\min_j\max_{|u|=1}m_j(u)\) allows the best direction to change between
pieces. Even strict positivity of this quantity leaves the common
direction undecided. Conversely the estimates themselves may vary
from chart to chart: if approximate linear maps \(A_j(q)\) in local
frames obey \(\|L_q-A_j(q)\|\le\epsilon_j\) and
\(\min_{C_j}|A_j(q)u|>\epsilon_j\) for the same \(u\), the triangle
inequality proves the required positivity.

\subsection{What a handle presentation supplies}

A handle presentation can be chosen on each individual smooth \(M\).
For example, Gay and Kirby's
trisection theorem \cite[Theorem~4]{GayKirby2016Trisecting}
gives an oriented closed connected smooth four-manifold a balanced
trisection of some type \((g_{\rm tr},k)\), with handle numbers
\[
       (n_0,n_1,n_2,n_3,n_4)=(1,k,g_{\rm tr}-k,k,1).
\]
On a homotopy sphere \(\chi=2\), so \(2+g_{\rm tr}-3k=2\).
One can thus start with counts \((1,k,2k,k,1)\) for some \(k\).
The attaching data, including the embeddings and framings of the
attaching solid tori for the two-handles, are still part of the input.
The identity of handle counts does not cancel these handles.

Give each \(D^i\times D^{4-i}\) product charts on its interior.
Near an attachment use charts on the attaching region and extend by a
signed collar coordinate, smoothing the product corners. If a local
attaching embedding is \(a\), the overlap in product collars is
represented, up to the choice of boundary coordinates, by
\((z,t)\mapsto(a(z),t)\) using a common signed normal coordinate.
Restricting such open sets to smaller compact pieces gives the cover
in Proposition~\ref{prop:cover}. A metric can likewise be chosen
using a partition of unity on this atlas, since positive-definite
quadratic forms form a convex cone. This supplies a conformal class;
the choice still affects \(L_q\) throughout \(X\).

The same construction works separately for every presented smooth
type. No classification of smooth structures is required as input.
If instead one transports the stereographic atlas by a homeomorphism
\(M\to S^4\), compatibility of the transported atlas with the given
smooth structure is precisely a smooth recognition question. The
handle construction respects the given smoothing from the outset.
Its transition functions organize the search for estimates; they
carry no lower bound on \(m(u)\) by themselves.

\subsection{Evaluation and generalized isothermal coordinates}

Fix a \(g(p)\)-orthonormal basis \(u_i\) of \(V\) and write
\(\omega_i=\omega_{u_i}\). The Gram matrix
\[
      Q_{ij}\,dV_g=\omega_i\wedge\omega_j,\qquad
      u^{\mathsf T}Q(q)u=|L_qu|_g^2
\]
is positive semidefinite. Its rank is the rank of \(L_q\). Where
\(L_qu\ne0\), the map \(u\mapsto L_qu/|L_qu|_g\) selects a point
of the twistor sphere \(S(\Lambda^+_q)\). If \(L_q\) is invertible,
normalization of \(L_q^{-1}\) gives the inverse map on spheres. At
lower rank the image lies in the unit sphere of a proper subspace.
Thus one global zero-free direction gives a smooth section of the
twistor sphere on \(X\); a chart covering the entire direction sphere
at \(q\) requires rank three there.

The additional condition for that whole-fibre chart to be conformal for the round sphere metrics is
\(Q(q)=\lambda(q)I\) with \(\lambda(q)>0\). Indeed, take orthonormal
right singular vectors of \(L_q\) with singular values \(s_1,s_2,s_3\).
At the first singular vector the derivative on the sphere stretches
the two orthogonal tangent directions by \(s_2/s_1\) and \(s_3/s_1\).
At the second it stretches them by \(s_1/s_2\) and \(s_3/s_2\).
Conformality forces all three values equal, and this condition suffices.

There is a differential consequence if scalar Gram holds throughout
an open set. The metric \(\widehat g=\sqrt{\lambda}\,g\) makes the
closed triple orthonormal. Locally orient the triple positively by a
constant sign change and set \(J_i b=\sqrt2\,*_{\widehat g}(b\wedge
\omega_i)\) on one-forms. In a standard self-dual frame
\(J_iJ_j=-\delta_{ij}+\sum_k\varepsilon_{ijk}J_k\). Write
\(\nabla\omega_i=\sum_j B_{ij}\otimes\omega_j\), \(B_{ij}=-B_{ji}\),
and put \((b_1,b_2,b_3)=(B_{23},B_{31},B_{12})\).
Alternating the covariant derivatives and using \(d\omega_i=0\)
gives \(J_2 b_3=J_3 b_2\) and its cyclic permutations. Thus
\(b_2=-J_3b_1\) and \(b_3=J_2b_1\); substituting in the first
equation gives \(2b_1=0\). The triple is parallel. This explains
the strength of requiring a conformal chart on every fibre.
Recognition only used a single coefficient line.

In the holomorphic correspondence-space formulation
motivated by the author's earlier work \cite{BakerCurved2025,BakerIsothermal2026},
a particularly strong generalized chart is a local biholomorphism
from \(U\times\mathbb{CP}^1\), defined on each entire \(\mathbb{CP}^1\) in its domain,
with base image in one holomorphic coordinate neighborhood, carrying
the flat \(\alpha\)-plane distribution to the curved distribution.
For each fixed point of \(U\) the projected coordinate functions are
holomorphic on \(\mathbb{CP}^1\), hence constant. The projection
therefore descends to a base map with invertible derivative. The
fibre maps are automorphisms of \(\mathbb{CP}^1\), so that derivative
preserves the full family of \(\alpha\)-planes and hence their union,
the null cone. The base map is conformal and its image is locally
conformally flat. A weaker construction can use separate patches in
the direction fibre. It would have to prove that the local data glue
and recover the intended curved geometry. An ordinary smooth handle
atlas and the real evaluation family do not supply that holomorphic
descent theorem.

\section{Smooth variation and zeros}\label{sec:pole-zeros}

\subsection{Moving the pole}

\begin{proposition}\label{prop:pole-openness}
On a fixed oriented smooth homotopy four-sphere, pole forms depend
continuously in every spatial derivative on compact sets away from the
moving pole. Along smooth finite-dimensional families of conformal
classes, poles and coefficients the dependence is smooth there. In the
bundle of nonzero coefficients over pairs \((c,p)\), the triples
\((c,p,\eta)\) with a zero-free pole form constitute an open set for
the smooth topology on conformal classes. The classes admitting such
a triple form an open set as well.
\end{proposition}

\begin{proof}
Fix a positive volume density to select a unique representative of
each class. Use nearby oriented orthonormal frames to trivialize
\(\Lambda^-_p\) and the framed exponential maps to compare nearby
punctures. More precisely, on a common small normal ball the composition
of the new framed exponential with the inverse reference exponential
sends the reference pole to the new pole. It extends to a diffeomorphism
of \(M\) by cutting off the displacement from the identity in a fixed
chart. For nearby data the displacement has derivative of norm less than
one, so the cut-off map has a smooth inverse by contraction and local
inversion.

Pull back to the reference punctured manifold, use the common inversion,
and multiply by the same factor \(|y|^{-4}\) on the end. In the framed
normal coordinates the compact metric is \(\delta_y+O(|y|^2)\), with
constant zeroth and first Taylor coefficients. Taylor's formula and
differentiation of the inversion give \(h=\delta_x+O_\infty(r^{-2})\)
for the blown-up metric, continuously in the corresponding symbol
bounds and smoothly along smooth finite-dimensional families. The
anti-self-dual target bundles of \(D=(d_h^*,\sqrt2\,d_h^-)\) can be
identified by projection onto the reference target. Take
\(\delta=-1/2\) in the weighted inverse used in A1. On a dyadic
annulus, rescaling to unit size bounds the operator difference
\(D-D_{\rm ref}\) by the differences of its rescaled coefficients in
\(C^k\), uniformly over the annuli. This gives continuity as a map
\(W^{k+1,2}_\delta\to W^{k,2}_{\delta-1}\). The inverse varies by the
convergent geometric series about the global reference inverse \(D_{\rm ref}^{-1}\). The same coefficient
estimates for parameter derivatives give smooth dependence of the
inverse in smooth families.

Let \(A^0=\chi\,\iota_x\beta/2\), where \(\chi=1\) far out and
\(\beta\) is the end two-form corresponding to \(\eta\). The wrong
projection \(d_h^-A^0\) is \(O_\infty(r^{-2})\). Its weighted target
norm on the annulus of radius \(R\) is \(O(R^{-1/2})\), so the right
side \((0,-\sqrt2\,d_h^-A^0)\) of the correction equation is continuous
in the target and smooth along smooth families. A1 identifies \(d(A^0+a)\), with \(a\)
the resulting correction, as the prescribed pole form. For \(k>2\),
Sobolev embedding on rescaled annuli gives
\[
 \sup_{A_R}|da|\le C R^{\delta-1}\|a\|_{W^{k+1,2}_\delta}.
\]
Thus \(\omega=\beta+O(r^{-3/2})\) uniformly for nearby data. Increasing
\(k\) gives the asserted dependence on compact sets. If the base form
is zero-free, its coefficient stays bounded away from zero nearby,
so one exterior radius suffices for the end. On the compact core the
base form has a positive minimum norm, which persists by uniform
convergence. Projection onto the class coordinate is open by local
triviality of the coefficient bundle.
\end{proof}

\subsection{Realization in a fixed inverted chart}

At a zero \(z\) of a two-form, the derivative
\(D_z\omega\colon T_zX\to\Lambda^2T_z^*X\) is intrinsic because the
terms from differentiating a local frame are multiplied by the vanishing
coefficients. If \(\omega\) is self-dual, the derivative takes values
in \(\Lambda_z^+\) by using a frame of that bundle. Thus a nonzero pole
form is positive symplectic off its zero set and is near-symplectic
precisely when this derivative has rank three at every zero, equivalently
when the section of \(\Lambda^+\) is transverse to the zero section.
Here the intrinsic near-symplectic condition for a closed form is
positive wedge-square away from its zeros and a wedge-positive
three-dimensional derivative image at each zero
\cite[Definition~1.1.1 and Proposition~1.1.4]{Perutz2005SurfaceFibrations}.

\begin{proposition}\label{prop:pole-marked-realization}
Let \(M\) be a closed oriented smooth four-manifold with a fixed
positive chart \(y\) centered at \(p\), and prescribe \(x=I(y)\) by
Euclidean inversion. Let \(\omega\) be near-symplectic or positive
symplectic on \(X=M\setminus\{p\}\), with \(\omega=\beta\)
sufficiently far out in \(x\) for a nonzero positive constant
two-form \(\beta\). A smooth conformal class on \(M\) makes this
\(\omega\) self-dual if and only if \(\beta\) is Euclidean self-dual
for the inherited orientation of \(X\). Its conformal value at
\(p\) is then \([\delta_y]\), and a compatible metric can be chosen
equal to \(\delta_y\) near \(p\).
\end{proposition}

\begin{proof}
The inherited orientation in \(x\) is opposite to the numerical
orientation since \(DI(\rho\nu)=\rho^{-2}R_\nu\), with \(R_\nu\)
the reflection normal to the unit vector \(\nu\). On each ray,
\(\rho^4\omega=R_\nu^*\beta\) in the \(y\) coframe. Continuity of
the compactified Hodge star forces all these reflected values into
the positive three-plane at \(p\). In an orthonormal coframe positively
oriented for the inherited end orientation, write
\(\beta=a e^{12}+b e^{34}\) with \(a,b>0\). Reflections normal to
\((e_1+e_3)/\sqrt2\) and \((e_2+e_4)/\sqrt2\) give
\(a e^{23}-b e^{14}\) and \(b e^{23}-a e^{14}\). If \(a\ne b\)
their span contains the nonzero square-zero form \(e^{23}\),
contradicting positivity. Hence \(a=b\). The reflected values are
then Euclidean self-dual in \(y\). The first reflection above together
with reflections normal to \(e_1\) and \((e_1+e_4)/\sqrt2\) gives
three independent such values, determining the plane at \(p\).

For sufficiency, use the correspondence between smooth wedge-positive
three-plane fields and Riemannian conformal classes
\cite[\S7.2, p.~1106]{AurouxDonaldsonKatzarkov2005SingularPencils}.
In a background metric splitting \(E_+\oplus E_-\), each positive
three-plane is the graph of a map \(B:E_+\to E_-\) with \(\|B\|<1\).
Containing the form imposes the affine equation
\(B\omega_+=\omega_-\). At a nonzero point the rank-one map
\(v\mapsto\langle v,\omega_+\rangle\omega_-/|\omega_+|^2\) works.
At a zero, first extend the derivative image to a local positive
plane field and use its splitting. The positive component of the form
is a submersion, so in local coordinates \((s,u)\) the components
are \((u,b(s,u))\). Wedge-positivity off the zeros forces
\(b(s,0)=0\). Hence
\[
 b(s,u)=\left(\int_0^1 \partial_u b(s,\lambda u)\,d\lambda\right)u.
\]
The operator in parentheses vanishes at the original zero by the choice
of splitting. It has norm less than one on a smaller neighbourhood.
The zero set there is \(u=0\), so the compact zero set is a finite
union of embedded circles when nonempty.

Use the Euclidean positive plane on the end. Express all local graphs
in a common background splitting and average with a partition of unity
whose end term equals one sufficiently far out. Convexity preserves
positivity and the affine equation. Choose the metric representative equal to \(\delta_x\) there. Since
\(I^*\delta_x=|y|^{-4}\delta_y\), the conformal rescaling by \(|y|^4\) near
\(p\) gives the compatible extension \(\delta_y\). On a homotopy sphere, \(y\) is now a normal chart and
A1 identifies the given form with the pole form of the specified
leading coefficient.
\end{proof}

The chart germ at \(p\), its Euclidean inversion and the end
coefficient are part of this realization. A linear change \(x'=Lx\)
induces the homogeneous map \(y'=|y|^2Ly/|Ly|^2\). Differentiability
at the origin would make it linear, so composing with \(L^{-1}\)
would give a linear map preserving each line and hence a scalar.
Thus the smooth compactifying changes of this kind are similarities.
For a general positive \(\beta\), one can instead modify the form
on the distant end. Choose \(L\in GL^+(4,\mathbb R)\) with
\(L^*\beta\) Euclidean self-dual. The flow generated by cutting off
\(\dot L_tL_t^{-1}x\) along a path from the identity to \(L\) fixes a
core containing the zeros. It exists for the full time interval by
at-most-linear growth. The least singular values of \(L_t\) are bounded
away from zero, so trajectories starting far enough out follow \(L_t x\).
Pullback gives the required end in the original chart. In the
other direction a pole form with transverse zeros has
\(\omega=\beta+da\) on the end by A1. Replacing \(a\) by
\(\chi_R a\) as in the recognition proof preserves positivity since
\(\chi_R da+d\chi_R\wedge a=O(R^{-2})\) on the transition annulus.
This gives an exactly constant end and preserves all its zeros.

\subsection{A first degenerate zero on the standard sphere}

\begin{proposition}\label{prop:pole-degenerate-first-zero}
On the standard smooth \(S^4\) there is a smooth path of metrics
with one fixed flat germ at a marked pole and one fixed nonzero
leading coefficient for which a zero first appears at a positive
parameter. Every zero at that parameter has intrinsic derivative of
rank at most two.
\end{proposition}

\begin{proof}
Choose \(y\) so its inverse \(x\) extends to global stereographic
coordinates on the punctured sphere. Identify the Euclidean \(x\)-space
with \(\mathbb C^2\) by an isometry respecting the inherited
orientation, and take \(\omega_0=\beta\) to be the constant K\"ahler
form. Taubes's local construction with zero initial circles gives a
closed form \(\omega_1\) with one zero circle and positive square
elsewhere, equal to \(\beta\) outside a ball
\cite[Proposition~2.1 and Part~5, following (2.16)]{Taubes2006VanishingCircles}.
Near the circle \(w=0,\ |z|^2=\varepsilon\), its coefficients in
the Euclidean self-dual plane are, up to constant invertible changes
of basis, \(-\varepsilon+|z|^2-|w|^2\) and the two real components
of \(zw-R\bar w\)
\cite[(2.2)--(2.5)]{Taubes2006VanishingCircles}.
The first coefficient has nonzero radial \(z\)-derivative there. The
last two have derivative \(v\mapsto zv-R\bar v\) in the \(w\)-plane
with real determinant \(\varepsilon-R^2\ne0\). Thus the derivative
image is the full Euclidean self-dual plane, as required by
Proposition~\ref{prop:pole-marked-realization}.

Choose compatible metrics \(g_0,g_1\) that both equal \(\delta_y\)
on a common neighbourhood of \(p\). The path
\(g_t=(1-t)g_0+tg_1\) preserves this germ, hence the chart and the
coefficient represented by \(\beta\). By A1 the endpoint pole forms
are exactly \(\omega_0,\omega_1\). By Proposition~\ref{prop:pole-openness}
and a finite parameter cover, a single compact core contains all zeros
of the family. The zero locus in the product of the interval with
this core is closed and compact. Its projection therefore has a
least element \(t_*>0\), since the first endpoint has no zero.

Near any zero at \(t_*\), a smooth frame of the family of self-dual
bundles represents the pole form by three functions. If the spatial
derivative had rank three, restricting to a three-dimensional slice
would make that derivative invertible. The implicit function theorem
would then give zeros on this slice for earlier nearby parameters,
contradicting minimality. If needed, the affine metric path extends
slightly past its endpoint by compactness and positive definiteness.
\end{proof}

\section{Supporting evidence on the standard sphere}
\label{sec:evidence}

\paragraph{An exact family.}
A model with a nonconstant metric illustrates how little symmetry of the
Gram matrix is needed for recognition. On \(\mathbb R^4\) write \(x=rq\),
\(q\in S^3\), and take a unit round coframe
\(\sigma_i=-(q^{-1}dq)_i\), so \(d\sigma_i=2\sigma_j\wedge\sigma_k\)
for cyclic \(i,j,k\). Orient by \(dr\wedge\sigma_1\wedge\sigma_2
\wedge\sigma_3\). For \(g=dr^2+\sum a_i^2\sigma_i^2\) put
\(\omega_i=d(f_i\sigma_i)\). In the orthonormal coframe
\((dr,a_i\sigma_i)\), self-duality reduces to the scalar equations
\[
       \frac{f_i'}{a_i}=\frac{2f_i}{a_j a_k}.
\]
Take \(t=r^2\), \(A=1+\varepsilon t/(1+t)^2\), \(a_1=rA\), and
\(a_2=a_3=r\). Then \(A>0\) exactly for \(\varepsilon>-4\) on the
whole radial interval. The solutions normalized by
\(f_i(r)/r^2\to1\) at infinity are \(f_i=r^2F_i(t)\), with
\[
 F_1(t)=e^{-\varepsilon/(1+t)},\qquad
 F_2(t)=F_3(t)=
 \exp\!\left(\varepsilon\int_t^\infty
          \frac{ds}{(1+s)^2+\varepsilon s}\right).
\]
Indeed the logarithmic derivatives of \(f_1,f_2,f_3\) in \(r\)
are \(2A/r,2/(rA),2/(rA)\), respectively. The \(F_i\) are smooth
and positive at zero. With \(\alpha_i=r^2\sigma_i\), a one-form
linear in Cartesian coordinates, \(d(F_i\alpha_i)\) at the origin
is \(F_i(0)d\alpha_i\); the three \(d\alpha_i\) form a constant
self-dual basis. Away from the origin the three solutions have
positive, diagonal Gram entries \(8f_i^2/(a_j a_k)^2\).

For smoothness of the metric at the two poles, its Cartesian expression
is \(g_{\rm euc}+[2\varepsilon/(1+t)^2+
\varepsilon^2 t/(1+t)^4]\alpha_1^2\), with
\(\alpha_1^2=\alpha_1\otimes\alpha_1\). Inversion preserves
\(t/(1+t)^2\); multiplication by \(4/(1+t)^2\) gives a smooth
representative on the compact sphere in both charts.
Since \(F_i=1+O(r^{-2})\) and \(dF_i=O(r^{-3})\) at infinity, the
three forms tend to \(d\alpha_i\) in this end chart. Inversion gives
quartic growth, so each belongs to the pole space. Their independent
end limits and A1 show that they form a basis. Thus every nonzero coefficient works at this fixed
pole on the full \(\varepsilon>-4\) family. This conclusion follows
from the differential equations. The sampled curvature integrals
reported in the proposal \cite{Baker2026PoleFormsProposal} are
numerical diagnostics.

\paragraph{A certified compactly supported perturbation.}
The second master studies a different class of metrics in an oriented
stereographic chart \((x_0,y)\). For its specified smooth translated
profile \(F\ge0\) the asymptotically Euclidean representative is
\[
     g_t=\left(\frac{1+tF}{1-tF}\right)^2dx_0^2+|dy|^2 .
\]
The profile, mollification and end coefficient basis are fixed in that
experiment. In flat chiral frames the graph of physical self-duality is
\(\omega^-=tF\omega^+\). The closed graph of the flat transfer then gives
\(\omega^+=u+T_{\rm flat}\omega^-\), with the constant identified with
the prescribed unit coefficient \(u\). This exhibits the same linear
transfer as in \eqref{eq:transfer} on the Euclidean end model. A
Fourier estimate and an independently replayed finite rational
certificate give \(|\omega_{t,u}(x)|_{g_t}\ge1/100\) for every \(x\),
every unit \(u\), and \(t\in[-3/4,0]\) for profile A.
At the endpoint \(t=-3/4\), separate enclosures on that class give
\(Y>2.21\) and \(0.00590\le E_+/(4\pi^2)\le1.00731\).
The pole estimate is uniform along the indicated segment. The quoted
curvature and Yamabe bounds concern the endpoint. The full argument
and the integer certificate are in the second master's
\emph{A uniform lower bound for the pole triple on profile A}
(27 November 2026). Neither main theorem uses a numerical estimate.

\paragraph{Draft provenance.}
Prepared for Edward B. Baker III with OpenAI Codex assistance on
27 November 2026, from the investigation's written proofs and independent
readings on repository base \texttt{6a422ce}. The exposition and
independent review used bounded \texttt{gpt-6.2-sol} helpers requested
at ultra effort; the exact serving variant and effort of the lead were
unavailable. The companion computations have the limited scope stated
above.

\section{Outstanding challenges}
\label{sec:open}

\paragraph{Selecting a direction on an arbitrary smooth type.}
For a fixed \((c,p)\), Proposition~\ref{prop:cover} describes the
entire obstruction as a compact subset of one projective plane.
The goal is to choose \(c,p\) so that \(\mathcal B(c,p)\) omits a projective point
\([u]\). An atlas from handles gives explicit domains on which to
estimate the same global solution. The missing step is an estimate on
the actual attaching regions and the interiors that leaves at least one
common direction. Its constants may depend on the presentation. A
bound on all three singular values everywhere would be sufficient;
a lower bound for just \(|L_qu|\) along the chosen \(u\) suffices.
Openness continues a witness once one is known. It supplies no
starting witness on an unknown smooth type.

\paragraph{What happens at a zero.}
A transverse zero set consists of circles; the first-zero argument of
Proposition~\ref{prop:pole-degenerate-first-zero} shows why one must
justify transversality when varying the original smooth data.
For a prepared single transverse twisted circle in the companion
positive cap construction, there exists a finite-energy
pseudoholomorphic current with one irreducible component of
multiplicity one.\footnote{The cap, chamber and period arguments are in
the companion investigation \emph{Lines and broken fibrations},
\emph{Relative degree on the marked projective cap} and \emph{A period test for components in the punctured
core} (27 November 2026). The marked line and the transverse circle
are essential hypotheses.} Its intersection degree with the
marked line is zero, while its algebraic meridional degree about the
circle has absolute value one. These are different pairings. Even
a unit meridional degree leaves open whether the closure is a
smooth embedded half-sheet all along the circle. To use such a
boundary for removal one must also control the interior of a
spanning surface and construct a deformation compatible with the
marked cap. In particular a local operation on the prepared form
would need to return a zero-free closed form on the specified smooth
compactification. The fixed-end criterion in
Proposition~\ref{prop:pole-marked-realization} describes that last
compatibility requirement.

\paragraph{Connecting curvature to the pole criterion.}
The first master \cite{Baker2026ChiralMaster} studies the band
\(Y(c)>0\), \(E_+(c)<8\pi^2\), with
\(E_+=\int_M |W^+|_{\rm HS}^2\,dV\) in the curvature-operator
normalization and \(Y(c)\) the Yamabe constant. Even on standard
\(S^4\) it is open whether \emph{each} class in this band has
\emph{some} pole and a zero-free direction. A failure at one
prescribed pole or a zero for one prescribed direction would leave
that question open. To obtain recognition on an unknown smooth type
by this route one would need the implication on that type; proving
entry into the band is a further, separate problem.

\paragraph{Coordinates and recognition.}
The stronger coordinate problem is to construct generalized charts
on patches of the correspondence space whose overlap data recover the
geometry and extend over the required direction fibres. A solution
would need hypotheses ensuring that descent and estimates connecting
it to the smooth Riemannian data. The pole criterion already has a
global closed form for each leading coefficient. Its remaining
question is exactly whether on every smooth type one of those forms
can be made symplectic everywhere on the punctured manifold.

\bibliographystyle{unsrturl}
\bibliography{references}
\end{document}
